\scanpage{79}\mspage{7\deleted{\unsure{7}}8}
At $9_{\Sigma}$ (critère de compatibilité par les lieux, pour $\Sigma$-ateliers) Soient $X,Y\in M$, \deleted{tels que} $\deleted{\unsure{K|X}}K=X\cap Y$ ($=\operatorname{Inf}^{\leq}(\unsure{x|X},\unsure{Y|y})=\operatorname{Sup}_{Z\in\widetilde X\cap\widetilde Y} Z$), de sorte que $\widetilde K=\widetilde X\cap\widetilde Y$. Pour que $x\deleted{\unsure{\text{et}|\parallel}}\mathrel{>\mkern-14mu<}y$, \deleted{soient disjoints} \inserted{\unsure{on|en}} il \deleted{\unsure{(}}f. et \deleted{\unsure{)}}s. que
\[
\begin{aligned}
\text{(i)}\quad &\forall\,x\in\operatorname{omb}(X)\setminus\operatorname{omb}(K),\quad y\in\unsure{\operatorname{om}|\operatorname{omb}}(Y)\setminus\operatorname{omb}(K),\\
&\text{on a}\quad x\parallel y\quad\text{(ou encore }x\mathrel{>\mkern-14mu<}y\text{)}\\
\text{(ii)}\quad &\widetilde K\in\Sigma
\end{aligned}
\]
Cela implique donc

\uline{Corollaire 1} (ancien At $9''$) Si on a (i) et (ii$'$) $K\in M$ (\unsure{i.e|i.e.} $\widetilde K$ a un plus grand élément), alors $X\mathrel{>\mkern-14mu<}Y$.

En effet, \unsure{.} $K\in M\Longrightarrow\widetilde K\in\Sigma$.

\uline{Corollaire 2} (ancien At $9'$) Supposons que $\Sigma\supset$ \deleted{\unsure{cat.}}catégorie des \unsure{ens|ens.} ordonnés isomorphes à un $\widetilde F$, $F\in\unsure{\mathcal F|\mathscr F}$ (p.ex. $\Sigma={}$cat. \unsure{de|des} tous les ens. ordonnés) alors $X\mathrel{>\mkern-14mu<}Y$ sss on a (i).

% slanted
\margin{\unsure{et|en} \unsure{prenant} l'un ou l'autre cas : 1$^{\mathrm o}$) $\Sigma={}$ tous les \unsure{ens.|ens} ordonnés, \unsure{donc} (ii) est vide, 2$^{\mathrm o}$) $\Sigma={}$ ss-cat. pleine des ens \unsure{avec} plus grand élément. (de sorte que (ii) équivaut à $K\in M$.)}

\uline{Scholie} \deleted{On supposera que tout atelier} \inserted{Tous les ateliers auxquels j'ai songé} \deleted{est un} \inserted{sont des} $\Sigma$-ateliers pour $\Sigma$ convenable, \deleted{ce qui \deleted{\unsure{implique que}} \inserted{\deleted{\unsure{équivaut}}} à dire que c'est un $\Sigma$-atelier pour $\Sigma\unsure{.}=\Sigma_M={}$cat. des ens. ord. équivalents \unsure{à|:} un $\widetilde X$, $X\in M$. Cela signifie aussi que At $9_{\Sigma}$ est valable pour $\Sigma_M$, ou encore que l'on veut \uline{\unsure{préciser les univers}} de validité du cor. 1}
